\documentclass[a4paper,12pt]{article}

\usepackage[T1]{fontenc}
\usepackage[italian]{babel}
\usepackage{graphicx}
\usepackage{geometry}
\usepackage{tabularx}
\usepackage[table]{xcolor}
\usepackage[most]{tcolorbox}
\usepackage{fancyhdr}
\usepackage[hidelinks]{hyperref}

\newcommand{\groupmail}{alittlebyte.swe@gmail.com}
\newcommand{\grouplogo}{../../../.github/site-src/assets/logo_alittlebyte.png}
\newcommand{\groupsymbol}{../../../.github/site-src/assets/solo_logo.png}

\geometry{
    left=2.5cm,
    right=2.5cm,
    top=2.5cm,
    bottom=2.5cm,
    headheight=331pt,
    includeheadfoot
}

\pagestyle{fancy}
\fancyhf{}

\renewcommand{\headrulewidth}{0pt}
\fancyhead{}

\renewcommand{\footrulewidth}{0.4pt}
\fancyfoot[L]{\includegraphics[height=0.6cm]{\groupsymbol}\hspace{0.45em}aLittleByte}
\fancyfoot[C]{\thepage}
\fancyfoot[R]{Verbale 2026-05-06}

\fancypagestyle{firstpage}{
    \fancyhf{}
    \renewcommand{\headrulewidth}{0pt}
    \renewcommand{\footrulewidth}{0.4pt}
    \fancyhead[C]{%
        \makebox[\textwidth][c]{%
            \begin{tabular}{c}
                \includegraphics[height=10cm]{\grouplogo}\\[1ex]
                \small\texttt{\groupmail}
            \end{tabular}%
        }
    }
    \fancyfoot[L]{\includegraphics[height=0.6cm]{\groupsymbol}\hspace{0.45em}aLittleByte}
    \fancyfoot[C]{\thepage}
    \fancyfoot[R]{Verbale  2026-05-06}
}

\providecommand{\glslink}[2]{\href{https://alittlebyte-19.github.io/Documentazione/glossario.html\#gls-#1}{\underline{#2}\textsuperscript{\scriptsize G}}}

\begin{document}

\thispagestyle{firstpage}

\vspace*{0.5cm}

\begin{center}
    {\LARGE \textbf{Verbale Esterno - 2026-05-06}}
\end{center}

\vspace{1.5cm}

\begin{center}
    \renewcommand{\arraystretch}{1.3}
    \begin{tcolorbox}[
        width=0.92\textwidth,
        colback=gray!6,
        colframe=gray!45,
        boxrule=0.5pt,
        arc=2mm,
        boxsep=0pt,
        left=8pt, right=8pt, top=8pt, bottom=8pt
    ]
        \begin{tabularx}{\linewidth}{>{\bfseries}p{3.5cm} X}
            Gruppo      & aLittleByte (gruppo 19) \\
            Redattore   & Alex Oloni \\
            Verificatori& E. Bellet\\
            Destinatari & Prof. Tullio Vardanega, Prof. Riccardo Cardin \\
            Data        & 2026-05-06\\
        \end{tabularx}
    \end{tcolorbox}
\end{center}

\newpage

\newgeometry{left=2.5cm,right=2.5cm,top=2.5cm,bottom=2.5cm,includefoot}
\tableofcontents
\newpage

\section{Informazioni Generali}
\section*{Specifiche Documento}
\hrule
\vspace{1cm}

\begin{tabular}{>{\bfseries}p{5cm} | p{10cm}}
Luogo Riunione & Google Meet \\[6pt]

Orario (Inizio - Fine) & 15:00 -- 15:20 \\[6pt]

\glslink{redattore}{Redattore} &
A. Oloni\\[6pt]

\glslink{amministratore}{Amministratore} & 
D. Belluz\\[6pt]

\glslink{verificatore}{Verificatore} &
E. Bellet\\[6pt]

Partecipanti &
  A. Oloni\\ &
  A. Gastaldello\\ &
  E. Bellet\\ & 
  D. Belluz\\ &
  A. Maule\\[6pt]

  Partecipanti Esterni &
  Luca Iuzzolino\\ &
  GianPaolo Ferrarin\\ [6pt]

\end{tabular}


\section{Ordine del Giorno}
\begin{itemize}
    \item Discussione e conferma delle tecnologie per lo sviluppo del PoC.
    \item Chiarimenti sulla gestione delle credenziali e degli accessi AWS.
    \item Approvazione dei \glslink{mockup}{mockup} e dell'\glslink{analisi-dei-requisiti-adr}{analisi dei requisiti}.
\end{itemize}


\section{Attività svolte}
Di seguito vengono riportate le richieste effettuate dal gruppo e le risposte fornite dal referente aziendale.

\subsection{Scelte Tecnologiche per il PoC}
Il gruppo ha proposto l'utilizzo di \textbf{\glslink{laravel}{Laravel}} e \textbf{\glslink{filament}{Filament}} per lo sviluppo del frontend e del backend. Il \glslink{committente-proponente}{proponente} ha confermato che non vi sono vincoli tecnologici e ha approvato pienamente la scelta, sottolineando che l'azienda dispone già di diversi progetti interni basati su questo stack.

\subsection{Gestione Account e Credenziali AWS}
È stata discussa la modalità di accesso alle risorse cloud. L'azienda fornirà al gruppo un account dedicato all'interno della propria organization AWS. Questo permetterà l'utilizzo di \glslink{amazon-bedrock}{Amazon Bedrock} e degli altri strumenti di \glslink{intelligenza-artificiale-ai}{intelligenza artificiale} necessari. Eventuali necessità di permessi aggiuntivi potranno essere gestite su richiesta.

\subsection{Approvazione Analisi dei Requisiti e Mockup}
Il \glslink{committente-proponente}{proponente} ha confermato di aver visionato l'\glslink{analisi-dei-requisiti-adr}{analisi dei requisiti} inviata tramite i canali di comunicazione informali. L'azienda ha espresso parere positivo, ritenendo il materiale prodotto esaustivo e corretto. Non è stata richiesta una firma formale sui documenti, poiché l'analisi è considerata uno strumento di gestione interna del gruppo.

\subsection{Presentazione dell'Interfaccia}
Durante l'incontro è stato mostrato il \glslink{mockup}{mockup} dell'applicativo, includendo la landing page e i due moduli core (\textit{IA Assistant} e \textit{Copilot}). Il referente ha apprezzato la chiarezza dell'interfaccia e la coerenza con le funzionalità attese (gestione \glslink{prompt}{prompt}, verifica manuale \glslink{ocr-riconoscimento-ottico-dei-caratteri}{OCR}, metriche).

\section{To-Do List}
\begin{table}[h]
    \centering
    \begin{tabular}{|l|p{5cm}|l|}
    \hline
        \rowcolor{lightgray}\textit{\textbf{Persona Incaricata}}  & \textbf{Task Assegnata} &\textbf{Link \glslink{issue}{Issue}} \\
    \hline
        \textit{Alex Oloni} & Redigere il verbale 2026-05-06 & \href{https://github.com/aLittleByte-19/Documentazione/issues/80}{\#80}\\
    \hline
    \end{tabular}
\end{table}

\end{document}
